\documentclass{apcalc}
\title{Unit 4 Practice Set: Contextual Applications of Differentiation}
\author{Paul Wang}
\date{November 2026}
\volume{Practice Set 4 of 10 (topics 4.1–4.7)}
\shorttitle{Unit 4 --- Practice Set}
\begin{document}
\maketitle
\begin{abstract}
Twelve exam-style multiple-choice questions on Unit 4, grouped by the sections of the Unit 4 Study Guide, and one nine-point free-response question on straight-line motion and L'Hôpital's rule. All items are no-calculator. Write each answer in the blank beside the number; the key on the last pages gives an answers-only strip, a reason and rule for every item, and the free-response scoring guide.
\end{abstract}
\scoringline{12 MC \quad\blank\ / 9 FR}{Exam pace: about 2 minutes per multiple-choice item (24 min) and 15 minutes for the free response.}
\section*{1\quad Derivatives in context}
\begin{setrule}R4.1–R4.3. Interpret with value, units and time. $v=x'$, $a=v'$; speed increases when $v$ and $a$ share a sign.\end{setrule}
\begin{q}{1}
$W(t)$ is the number of gallons of water in a tank $t$ hours after noon. Which is the best interpretation of $W'(4)=-6$?
\chone{Four hours after noon the tank holds 6 fewer gallons than it did at noon.}{At 4 p.m. the amount of water in the tank is decreasing at 6 gallons per hour.}{During the first 4 hours, water leaves the tank at an average rate of 6 gallons per hour.}{At 4 p.m. the rate at which water leaves the tank is decreasing by 6 gallons per hour per hour.}
\end{q}
\begin{q}{2}
A particle moves along the $x$-axis with position $x(t)=t^3-9t^2+24t$ for $t\ge0$. For which values of $t$ is the speed of the particle increasing?
\chtwo{$2<t<4$ only}{$t>3$ only}{$0<t<2$ and $3<t<4$}{$2<t<3$ and $t>4$}
\end{q}
\begin{q}{3}
A particle has position $x(t)=2t^3-15t^2+36t$. What is its acceleration the first time it is at rest?
\chfour{$-6$}{$0$}{$6$}{$28$}
\end{q}
\begin{q}{4}
$P(t)$ is the population of a town, $t$ years after 2020. Which is the best interpretation of $P''(5)=30$?
\chone{In 2025 the population increased by 30 people.}{In 2025 the population was 30.}{In 2025 the rate of population growth was increasing at 30 people per year per year.}{In 2025 the population was increasing at 30 people per year.}
\end{q}
\section*{2\quad Related rates}
\begin{setrule}R4.4–R4.5. Relate the quantities, differentiate with respect to $t$, then substitute the instant's values.\end{setrule}
\begin{q}{5}
The radius of a sphere increases at 0.5 cm per second. How fast is the volume increasing when the radius is 4 cm? ($V=\frac43\pi r^3$)
\chfour{$8\pi$ cm³/s}{$32\pi$ cm³/s}{$64\pi$ cm³/s}{$\frac{256}{3}\pi$ cm³/s}
\end{q}
\begin{q}{6}
A 13-foot ladder leans against a wall. The bottom slides away from the wall at 2 ft/s. How fast is the top sliding down the wall when the bottom is 5 ft from the wall?
\chfour{$\frac{5}{12}$ ft/s}{$\frac56$ ft/s}{$2$ ft/s}{$\frac{24}{5}$ ft/s}
\end{q}
\begin{q}{7}
Water is poured at 3 m³/min into a conical tank, vertex down, 10 m tall with top radius 5 m. How fast is the water level rising when the water is 6 m deep?
\chfour{$\dfrac{1}{48\pi}$ m/min}{$\dfrac{1}{9\pi}$ m/min}{$\dfrac{1}{3\pi}$ m/min}{$\dfrac{1}{\pi}$ m/min}
\end{q}
\section*{3\quad Local linearity and approximation}
\begin{setrule}R4.6–R4.7. $f(x)\approx f(a)+f'(a)(x-a)$; concave up means underestimate, concave down means overestimate.\end{setrule}
\begin{q}{8}
$f$ is differentiable with $f(2)=5$ and $f'(2)=-3$. The tangent-line approximation of $f(2.2)$ is
\chfour{$4.4$}{$4.7$}{$5.0$}{$5.6$}
\end{q}
\begin{q}{9}
$f(1)=2$, $f'(1)=3$ and $f''(x)=-x^2-1$ for all $x$. Using the line tangent at $x=1$, the approximation of $f(1.2)$ is
\chtwo{$2.6$, an underestimate}{$2.3$, an overestimate}{$3.6$, an overestimate}{$2.6$, an overestimate}
\end{q}
\section*{4\quad L'Hôpital's rule}
\begin{setrule}R4.8–R4.9. Confirm $\frac00$ or $\frac\infty\infty$, then $\lim\frac fg=\lim\frac{f'}{g'}$; rewrite $1^\infty$ and $0\cdot\infty$ first.\end{setrule}
\begin{q}{10}
$\displaystyle\lim_{x\to0}\frac{1-\cos 2x}{x^2}$ is
\chfour{$0$}{$1$}{$2$}{$4$}
\end{q}
\begin{q}{11}
Differentiable $f$ and $g$ satisfy $f(3)=g(3)=0$, $f'(3)=4$ and $g'(3)=-2$. What is $\displaystyle\lim_{x\to3}\frac{f(x)}{g(x)}$?
\chfour{$-2$}{$-\dfrac12$}{$0$}{nonexistent}
\end{q}
\begin{q}{12}
$\displaystyle\lim_{x\to\infty}\left(1+\frac2x\right)^{3x}$ is
\chfour{$1$}{$e^2$}{$e^5$}{$e^6$}
\end{q}
\clearpage\section*{Free response}
Show all work. Justify every conclusion with calculus, and give decimal answers correct to three places.
\begin{frq}{Question 1 --- Straight-line motion (9 points, no calculator)}
A particle moves along the $x$-axis so that its position at time $t\ge0$ is $x(t)=2t^3-9t^2+12t$.
\partlabel{a}Find the velocity and the acceleration of the particle at $t=\frac54$. Is the speed of the particle increasing or decreasing at $t=\frac54$? Give a reason for your answer.\workspace{1.6in}
\partlabel{b}For $0<t<3$, at what times does the particle change direction? Justify your answer.\workspace{1.6in}
\partlabel{c}Find the total distance the particle travels over $0\le t\le3$.\workspace{1.6in}
\partlabel{d}Find $\displaystyle\lim_{t\to1}\frac{x(t)-5}{(t-1)^2}$.\workspace{1.6in}
\end{frq}
\clearpage\section*{Answer key}
\keystrip{\keyitem{1}{B}\keyitem{2}{D}\keyitem{3}{A}\keyitem{4}{C}\keyitem{5}{B}\keyitem{6}{B}\keyitem{7}{C}\keyitem{8}{A}\keyitem{9}{D}\keyitem{10}{C}\keyitem{11}{A}\keyitem{12}{D}}
\vspace{4pt}
\begin{keytable}
\keysection{1. Derivatives in context}
1 & B & A derivative is an instantaneous rate: value, units (gallons per hour), and the time. & R4.1\\
2 & D & $v=3(t-2)(t-4)$, $a=6t-18$; same signs on $(2,3)$ (both negative) and $t>4$ (both positive). & R4.2\\
3 & A & $v=6(t-2)(t-3)$; first rest at $t=2$; $a(2)=24-30=-6$. & R4.2\\
4 & C & $P''$ is the rate of change of the growth rate: people per year per year. & R4.3\\
\keysection{2. Related rates}
5 & B & $\frac{dV}{dt}=4\pi r^2\frac{dr}{dt}=4\pi(16)(0.5)$. & R4.5\\
6 & B & $xx'+yy'=0$ with $x=5$, $y=12$, $x'=2$: $y'=-\frac{10}{12}$. & R4.4\\
7 & C & $r=\frac h2$, $V=\frac{\pi h^3}{12}$, $3=\frac{\pi h^2}{4}h'=9\pi h'$. & R4.5\\
\keysection{3. Local linearity and approximation}
8 & A & $5+(-3)(0.2)=4.4$. & R4.6\\
9 & D & $2+3(0.2)=2.6$; $f''<0$, so the tangent lies above the curve. & R4.7\\
\keysection{4. L'Hôpital's rule}
10 & C & L'Hôpital twice: $\frac{2\sin2x}{2x}\to\frac{4\cos2x}{2}=2$. & R4.9\\
11 & A & $\frac00$, so the limit is $\frac{f'(3)}{g'(3)}=-2$. & R4.9\\
12 & D & $\ln y=3x\ln(1+\frac2x)\to6$. & R4.8\\
\end{keytable}
\section*{Free-response scoring guide}
\subsection*{Question 1: Straight-line motion (9 points)}
\textbf{(a)} $v(t)=6t^2-18t+12=6(t-1)(t-2)$, so $v(\frac54)=6(\frac14)(-\frac34)=-\frac98$. $a(t)=12t-18$, so $a(\frac54)=-3$. Both are negative, so the speed is increasing at $t=\frac54$.
\begin{itemize}[leftmargin=1.6em,itemsep=0pt]
\item 1 point: $v(\frac54)=-\frac98$
\item 1 point: $a(\frac54)=-3$
\item 1 point: speed increasing, because $v$ and $a$ are both negative
\end{itemize}
\textbf{(b)} $v=6(t-1)(t-2)$ changes from positive to negative at $t=1$ and from negative to positive at $t=2$, so the particle changes direction at $t=1$ and $t=2$.
\begin{itemize}[leftmargin=1.6em,itemsep=0pt]
\item 1 point: $t=1$ and $t=2$
\item 1 point: justification: $v$ changes sign at each
\end{itemize}
\textbf{(c)} $x(0)=0$, $x(1)=5$, $x(2)=4$, $x(3)=9$. Distance $=|5-0|+|4-5|+|9-4|=5+1+5=11$.
\begin{itemize}[leftmargin=1.6em,itemsep=0pt]
\item 1 point: positions at $t=0,1,2,3$: $0,5,4,9$
\item 1 point: total distance 11
\end{itemize}
\textbf{(d)} $x(1)-5=0$ and $(1-1)^2=0$. L'Hôpital: $\lim\frac{v(t)}{2(t-1)}$, again $\frac00$ since $v(1)=0$. Again: $\lim\frac{a(t)}{2}=\frac{a(1)}{2}=\frac{-6}{2}=-3$. (Or factor: $x(t)-5=(t-1)^2(2t-5)$.)
\begin{itemize}[leftmargin=1.6em,itemsep=0pt]
\item 1 point: shows the form is $\frac00$ and applies L'Hôpital's rule (or factors)
\item 1 point: limit $-3$
\end{itemize}
\end{document}